%!TEX root = ../../thesis.tex

\section{Overview}\label{sec:corpus-overview}

% NOTE: Keep the six sections of this chapter aligned with
% \cref{chap:input-actions} --- same headings, same order --- so the two formulations
% read as a diff.

% TODO: State the formulation: the agent operates the greybox fuzzing loop itself.
% It maintains a fixed-size corpus and each step is one composite action
% (retain, select, position, operator) applied in the order retain -> select ->
% mutate. Give the step loop as an algorithm, laid out like the one in
% \cref{sec:input-overview}.
%
% Running example for this chapter: cjson under the corpus environment
% (\cref{sec:target-cjson}), which is also the target \cref{sec:results-head-to-head}
% uses to compare the two formulations directly.
%
% Motivate the formulation from \cref{sec:fuzzing-cgf}: selection, mutation and
% retention are the three problems every coverage-guided greybox fuzzer solves with
% hand-tuned heuristics. Here all three are learned decisions. That, rather than the
% mutation operators, is the novel part --- \cite{bottinger_deep_2018} learned the
% mutation and left selection random and retention absent.
